% Source: https://artofproblemsolving.com/community/c3h2184638_interesting_problem
% Stable identity: aops-topic-2184638-q01
Multiplying the first equation by $c^{\dfrac{3}{2}}=c\sqrt c$ gives the second. Suppose, for a contradiction, that $a\sqrt a+b\sqrt b=c\sqrt c$ with positive integers $a,b,c$.

\textbf{Step 1: $ab$ and $bc$ are perfect squares.} Squaring,
\[a^3+2ab\sqrt{ab}+b^3=c^3,\qquad\text{so}\qquad\sqrt{ab}=\dfrac{c^3-a^3-b^3}{2ab}\]
is rational. The square root of a positive integer is rational only if the integer is a perfect square, so $ab$ is a perfect square. Similarly, squaring $c\sqrt c-b\sqrt b=a\sqrt a$ gives $c^3-2bc\sqrt{bc}+b^3=a^3$, so $\sqrt{bc}$ is rational and $bc$ is a perfect square.

\textbf{Step 2: common square-free part.} Write $b=dy^2$ where $d$ is square-free (the product of the primes that divide $b$ to an odd power). Since $ab=a\,d\,y^2$ is a perfect square, so is $ad$ (it equals $\left(\dfrac{\sqrt{ab}}{y}\right)^2$, a rational square that is an integer). For a prime $p$ dividing $d$, the exponent of $p$ in $ad$ is even, so its exponent in $a$ is odd; for a prime not dividing $d$, its exponent in $a$ is even. Hence $\dfrac {a}{d}$ is an integer with all exponents even: $a=dx^2$ for a positive integer $x$. In the same way, from $bc$ being a square, $c=dz^2$.

\textbf{Step 3: reduce to Fermat.} Now $a\sqrt a=dx^2\cdot x\sqrt d=d\sqrt d\,x^3$, and similarly for $b$ and $c$. The equation becomes
\[d\sqrt d\left(x^3+y^3\right)=d\sqrt d\,z^3,\qquad\text{so}\qquad x^3+y^3=z^3\]
with positive integers $x,y,z$. This contradicts Fermat's Last Theorem for exponent $3$. Hence no such $a,b,c$ exist. $\blacksquare$
